Five classifiers achieved perfect discrimination on the UCI CKD benchmark, then failed every stress-test criterion when applied to the MIMIC-IV demo cohort under deliberate distributional shift. Isotonic recalibration reduced internal ECE to near zero. On the demo cohort, calibration drift exceeded 0.67 for every model, conformal coverage fell from 0.967 or higher internally to 0.21--0.25, and no model scored above 4 out of 16 on the deployment readiness checklist.

The MIMIC cohort used here is an open-access 100-patient demonstration set, not a formally designed external validation dataset. The stress-test framing is intentional. Its purpose is to illustrate that internal calibration quality, however strong, does not guarantee the same reliability under prevalence shift and feature missingness. That point holds regardless of the demo cohort's limitations.

The eight-criterion framework introduced here offers one structured path for evaluating these requirements before clinical use. Each criterion has a defined threshold and a three-level scoring scheme applicable to any binary clinical prediction model. Formal external validation on a purpose-built cohort with properly measured features and an appropriate sample size is the necessary next step. The framework provides the evaluation structure. The full MIMIC-IV database provides the cohort.
